\section{Experiments}\label{sec:experiments}

\subsection{rMD17 aspirin main result}\label{sec:rmd17}

We evaluate the main molecular setting on the rMD17 aspirin trajectory dataset \cite{chmiela2017}, an MD trajectory of aspirin with 100k frames. We subsample every 10th frame to obtain 10k frames, and from these we sample 2000 training transitions and 200 evaluation transitions per seed, where a transition is a consecutive (t, t+1) frame pair. Train and evaluation transition starting frame indices are sampled disjointly per seed. The world model uses a flat encoder, latent transition predictor, and reward/done heads. All priors use latent dimension 16, batch size 32, 50 epochs, prior weight \(\lambda=0.1\), and the default learning rate. For aspirin, the bond graph contains 21 atoms, with covalent bonds as edges and graph Laplacian \(L=D-A\). We report latent rollout error \(\|\hat{z}_{t+H}-z_{t+H}\|_2^2\), where \(\hat{z}_{t+H}\) is produced by unrolling the latent transition model and \(z_{t+H}\) is encoded from the ground-truth observation. We evaluate \(H \in \{1,2,4,8,16\}\) over 10 seeds per prior.

Table~\ref{tab:rmd17} and Figure~\ref{fig:rmd17} show that the graph-spectral prior reduces rollout error across the full horizon range. At \(H=16\), the spectral prior achieves \(0.140 \pm 0.065\), compared with \(0.290 \pm 0.193\) for no prior and \(0.323 \pm 0.139\) for the Euclidean prior. This corresponds to a 51.5\% reduction relative to no prior and a 56.5\% reduction relative to Euclidean. The 95\% confidence intervals are also strongly separated in practical terms: the spectral interval \([0.103, 0.178]\) lies almost entirely below the no-prior interval \([0.176, 0.412]\), with only a narrow boundary overlap. Importantly, the effect is not restricted to long-horizon rollout. At \(H=1\), the spectral prior already reduces error from 0.101 to 0.037, a 63\% reduction. Across \(H=1\) through \(H=16\), the relative reduction remains in the 50--65\% range. This pattern suggests that the spectral prior improves latent representation quality itself, rather than only stabilizing accumulated transition errors. Figure~\ref{fig:rmd17} visualizes the same trend, with the spectral curve and confidence band consistently below the no-prior and Euclidean curves.

The spectral prior also reduces seed-to-seed variability. At \(H=16\), its standard deviation is 0.065, compared with 0.193 for no prior, a 3.0x reduction. At \(H=1\), the reduction is larger: 0.005 for spectral versus 0.052 for no prior, or 10.4x lower variability. Thus, the spectral prior not only achieves lower mean rollout error but also produces more reliable seed-to-seed performance. This matters for downstream uses of learned world models, including model-based reinforcement learning and planning, where consistency across random initializations is often as important as average performance.

At \(\lambda=0.1\), the Euclidean covariance prior performs slightly worse than the no-prior baseline at every horizon. For example, at \(H=16\), it obtains \(0.323 \pm 0.139\) versus \(0.290 \pm 0.193\) for no prior, or +11.4\%. The Euclidean prior is sensitive to \(\lambda\) choice; in Section~\ref{sec:weight_sweep} we find its optimum at \(\lambda=0.01\), where it achieves \(H=16\) error \(0.173 \pm 0.049\) with \(n=3\). Even at this tuned setting, the Euclidean prior does not reach the spectral prior's \(H=16\) error of \(0.140 \pm 0.065\) in Table~\ref{tab:rmd17}. We report all priors at the shared \(\lambda=0.1\) in this section for direct comparison; the weight sensitivity analysis in Section~\ref{sec:weight_sweep} confirms that the spectral prior maintains its advantage under per-prior weight tuning.
